\SourceLeaf{39}\Article{1}
Zur Aufstellung der Gleichgewichtsgleichung eines Systems aus beliebig vielen materiellen Punkten, deren Bewegungen beliebigen Bedingungen unterliegen, eignet sich am besten das Prinzip der virtuellen Bewegungen. Wir formulieren es wie folgt.

Das System bestehe aus materiellen Punkten $m,m',m''$ usw., in denen wir uns die mit denselben Buchstaben bezeichneten Massen konzentriert denken. Es sei $P$ eine der beschleunigenden Kräfte, die auf den Punkt $m$ wirken. Wir denken uns, dass das System eine beliebige unendlich kleine, mit seinen Bedingungen vereinbare Bewegung ausführt, also eine virtuelle Bewegung. Dabei sei $\mathrm{d}p$ die Bewegung des Punkts $m$ in der Projektion auf die Richtung der Kraft $P$, also die Bewegung multipliziert mit dem Kosinus des Winkels, den sie mit der Richtung der Kraft $P$ bildet. Schließlich sei $\Sigma P\mathrm{d}p$ die Summe aller derartigen Produkte für alle auf den Punkt $m$ wirkenden Kräfte. Ebenso bezeichne $P'$ allgemein die auf den Punkt $m'$ wirkenden Kräfte und $\mathrm{d}p'$ die auf deren jeweilige Richtung projizierten Bewegungen des Punkts $m'$, und entsprechend für die übrigen Punkte. Mit diesen Bezeichnungen besteht die Gleichgewichtsbedingung des Systems darin, dass die Summe
\[
m\Sigma P\mathrm{d}p+m'\Sigma P'\mathrm{d}p'+m''\Sigma P''\mathrm{d}p''+\ \text{etc.}
\]
für jede virtuelle Bewegung $=0$ wird, wie das Prinzip der virtuellen Bewegungen gewöhnlich formuliert wird. Genauer gesagt kann diese Summe für keine virtuelle Bewegung einen positiven Wert annehmen.

\Article{2}
Die hier zu betrachtenden Kräfte lassen sich in drei Klassen einteilen.

I. Die Schwerkraft. Man darf annehmen, dass ihre Stärke für alle Punkte dieselbe ist und ihre Richtungen parallel sind. Ihre Stärke bezeichnen wir mit $g$.

\SourceLeaf{40}II. Die Anziehungskräfte, welche die Punkte $m,m',m''$ usw. aufeinander ausüben. Es wird angenommen, dass die Stärke der Anziehung proportional zu einer Funktion der Entfernung ist, oder, anders ausgedrückt, dass sie gleich dem Produkt dieser mit dem Zeichen $f$ bezeichneten Funktion und der im anziehenden Punkt konzentrierten Masse ist.

III. Die Kräfte, durch welche die Punkte $m,m',m''$ usw. von beliebig vielen festen Punkten angezogen werden. Für diese Kräfte verwenden wir entsprechend das der Entfernung vorangestellte Zeichen $F$ und bezeichnen mit $M,M',M''$ usw. sowohl die festen Punkte als auch die Massen, die wir uns in ihnen konzentriert denken.

Bezeichnen wir nun den Abstand der beiden Punkte $m,m'$ mit $(m,m')$ und ebenso mit $(m,M)$ den Abstand der Punkte $m,M$ usw., ferner mit $z,z',z''$ usw. die Höhen der Punkte $m,m',m''$ usw. über einer beliebigen horizontalen Ebene $H$, so erhalten wir für den Ausdruck $\Sigma P\mathrm{d}p$ die folgenden Bestandteile:
\par\begingroup\small
\[
\begin{aligned}
&-g\mathrm{d}z\\
&-m'f(m,m')\mathrm{d}(m,m')-m''f(m,m'')\mathrm{d}(m,m'')-m'''f(m,m''')\mathrm{d}(m,m''')-\ \text{etc.}\\
&-MF(m,M)\mathrm{d}(m,M)-M'F(m,M')\mathrm{d}(m,M')-M''F(m,M'')\mathrm{d}(m,M'')-\ \text{etc.}
\end{aligned}
\]
\endgroup
wobei die Differenziale $\mathrm{d}(m,m'),\mathrm{d}(m,m'')$ usw. partiell sind, da sie sich nur auf die virtuelle Bewegung des Punkts $m$ beziehen.

Anstelle der Funktion $f$ wollen wir nun diejenige Funktion einführen, aus der sie durch Differenziation entsteht. Wir setzen also $-fx.\mathrm{d}x=\mathrm{d}\varphi x$ oder $\int fx.\mathrm{d}x=-\varphi x$. Die Integrationskonstante lässt sich nach Belieben wählen. Wenn man will und wenn es die Sache zulässt, kann man sie so bestimmen, dass $\varphi\infty=0$ gilt. In diesem Fall stellt $\varphi t$ das Integral $\int fx.\mathrm{d}x$ von $x=t$ bis $x=\infty$ dar. Ganz entsprechend wird anstelle der Funktion $F$ eine andere Funktion $\Phi$ eingeführt, für die $-Fx.\mathrm{d}x=\mathrm{d}\Phi x$ gilt. Damit wird der Ausdruck $\Sigma P\mathrm{d}p$ zu $=$
\[
\begin{aligned}
&-g\mathrm{d}z\\
&+m'\mathrm{d}\varphi(m,m')+m''\mathrm{d}\varphi(m,m'')+m'''\mathrm{d}\varphi(m,m''')+\ \text{etc.}\\
&+M\mathrm{d}\Phi(m,M)+M'\mathrm{d}\Phi(m,M')+M''\mathrm{d}\Phi(m,M'')+\ \text{etc.}
\end{aligned}
\]
Dabei ist zu beachten, dass die Differenziale in der zweiten Zeile partiell sind und sich nur auf die Bewegung des Punkts $m$ beziehen.

Offenbar hat aber jedes dieser partiellen Differenziale seine Ergänzung in einem anderen Ausdruck. So enthalten sowohl der Ausdruck $m\Sigma P\mathrm{d}p$ als auch der Ausdruck $m'\Sigma P'\mathrm{d}p'$ das partielle Differenzial $mm'\mathrm{d}\varphi(m,m')$. Im ersten bezieht es sich nur auf die Bewegung von $m$, im zweiten nur auf die von $m'$. Daraus geht hervor, dass die in Artikel 1 angegebene Sum\SourceLeaf{41}me tatsächlich ein vollständiges Differenzial ist, und zwar $=\mathrm{d}\Omega$, wenn wir setzen $\Omega=$
\[
\begin{aligned}
&-gmz-gm'z'-gm''z''-\ \text{etc.}\\
&+mm'\varphi(m,m')+mm''\varphi(m,m'')+mm'''\varphi(m,m''')+\ \text{etc.}\\
&\phantom{+mm'\varphi(m,m')}+m'm''\varphi(m',m'')+m'm'''\varphi(m',m''')+\ \text{etc.}\\
&\phantom{+mm'\varphi(m,m')+mm''\varphi(m,m'')}+m''m'''\varphi(m'',m''')+\ \text{etc.}\\
&\phantom{+mm'\varphi(m,m')+mm''\varphi(m,m'')+mm'''\varphi(m,m''')}+\ \text{etc.}\\
&+mM\Phi(m,M)+mM'\Phi(m,M')+mM''\Phi(m,M'')+\ \text{etc.}\\
&+m'M\Phi(m',M)+m'M'\Phi(m',M')+m'M''\Phi(m',M'')+\ \text{etc.}\\
&+m''M\Phi(m'',M)+m''M'\Phi(m'',M')+m''M''\Phi(m'',M'')+\ \text{etc.}\\
&+\ \text{etc.}
\end{aligned}
\]
Die Gleichgewichtsbedingung besteht also darin, dass der Wert der Funktion $\Omega$ durch keine virtuelle Bewegung einen positiven Zuwachs erhalten kann, oder, was dasselbe ist, \textit{dass $\Omega$ ein Maximum ist}.

Die Funktion $\Omega$ lässt sich auch folgendermaßen darstellen:
\[
\begin{aligned}
\Omega=\Sigma m\{&-gz+\tfrac12m'\varphi(m,m')+\tfrac12m''\varphi(m,m'')+\tfrac12m'''\varphi(m,m''')+\ \text{etc.}\\
&+M\varphi(m,M)+M'\Phi(m,M')+M''\Phi(m,M'')+\ \text{etc.}\}
\end{aligned}
\]
Das Zeichen $\Sigma$ steht für die Summe aus dem daneben geschriebenen Ausdruck und allen Ausdrücken, die daraus hervorgehen, wenn $m$ nacheinander mit $m',m'',m'''$ usw. vertauscht wird.

